% Môn: Vật lý
% Lớp: 11
% Chương: Chương III. Điện trường
% Bài: L11C3  Bài 18. Điện trường đều
% Số câu: 11
% App đọc file này trực tiếp — sửa rồi Commit trên GitHub, bấm Đồng bộ GitHub .tex

% L11C3  Bài 18. Điện trường đều

% ===== Câu 1 =====
% ID: L11C3B18-01-TL
% Mức: NB
\dangbt{Chưa có dạng}
\begin{ex}
% ID: L11C3B18-01-TL
% Mức: NB
Hãy cho ví dụ về ứng dụng thực tiễn tác dụng của điện trường đối với chuyển động của điện tích bay vào điện trường đều theo phương vuông góc với đường sức. $\nguon{ly11 kntt.pdf}$
\choiceTF

\loigiai{
- Máy lọc không khí sử dụng công nghệ ion âm sẽ phát ra các ion âm vào trong không khí. Điện trường đều của Trái Đất làm phân tán rộng chùm ion âm này và hướng chúng lên phía trên. Tác dụng này làm tăng khả năng để các ion âm kết hợp được với các hạt bụi mịn mang điện dương tức là tăng khả năng lọc bụi mịn.

- Trong dao động kí, điện trường đều của các bản lái tia có tác dụng điều chỉnh hướng đi của các tia điện tử (electron).
}
\end{ex}

% ===== Câu 2 =====
% ID: L11C3B18-01-TN
% Mức: NB
\begin{ex}
% ID: L11C3B18-01-TN
% Mức: NB
Cường độ điện trường đều giữa hai bản kim loại phẳng song song được nối với nguồn điện có hiệu điện thế $U$ sẽ giảm đi khi $\nguon{ly11 kntt.pdf}$
\choice
{tăng hiệu điện thế giữa hai bản phẳng.}
{\True tăng khoảng cách giữa hai bản phẳng.}
{tăng diện tích của hai bản phẳng.}
{giảm diện tích của hai bản phẳng.}
\loigiai{
Chọn B.

Công thức: $E = \dfrac{U}{d}$ nên cường độ điện trường đều giữa hai bản kim loại phẳng song song được nối với nguồn điện có hiệu điện thế $U$ sẽ giảm đi khi tăng khoảng cách giữa hai bản phẳng.
}
\end{ex}

% ===== Câu 3 =====
% ID: L11C3B18-02-TL
% Mức: TH
\begin{ex}
% ID: L11C3B18-02-TL
% Mức: TH
Trong cơ thể sống, có nhiều loại tế bào, màng tế bào có nhiệm vụ kiểm soát các chất và ion ra vào tế bào đảm bảo cho quá trình trao đổi chất và bảo vệ tế bào trước các tác nhân có hại của môi trường. Một tế bào có màng dày khoảng 8 $\cdot$ 10^{-9} m, mặt trong của màng tế bào mang điện tích âm, mặt ngoài mang điện tích dương. Hiệu điện thế giữa hai mặt này bằng $0,07\,\mathrm{V}$. Hãy tính cường độ điện trường trong màng tế bào trên. $\nguon{ly11 kntt.pdf}$
\choiceTF

\loigiai{
Áp dụng công thức $\text{E} = \dfrac{U}{\text{d}}$ ta tính được cường độ điện trường trong màng tế bào bằng: $E = \dfrac{0,07}{8\cdot 10^{-9}} = 8,75\cdot 10^{6}\text{V}/\text{m}$
}
\end{ex}

% ===== Câu 4 =====
% ID: L11C3B18-02-TN
% Mức: NB
\begin{ex}
% ID: L11C3B18-02-TN
% Mức: NB
Các đường sức điện trong điện trường đều $\nguon{ly11 kntt.pdf}$
\choice
{chỉ có phương là không đổi.}
{chỉ có chiều là không đổi.}
{\True là các đường thẳng song song cách đều.}
{là những đường thẳng đồng quy.}
\loigiai{
Chọn C.

Các đường sức điện trong điện trường đều là các đường thẳng song song cách đều.
}
\end{ex}

% ===== Câu 5 =====
% ID: L11C3B18-03-TL
% Mức: TH
\begin{ex}
% ID: L11C3B18-03-TL
% Mức: TH
Một ion âm có điện tích -3,2 $\cdot10^{-19}C$ đi vào trong màng tế bào ở câu 7. Hãy xác định xem ion âm sẽ bị đẩy ra khỏi tế bào hay đẩy vào trong tế bào và lực điện tác dụng lên ion âm bằng bao nhiêu. $\nguon{ly11 kntt.pdf}$
\choiceTF

\loigiai{
Điện trường trong màng tế bào sẽ hướng từ phía ngoài vào trong. Vì lực tác dụng lên ion âm ngược chiều với cường độ điện trường nên lực điện sẽ đẩy ion âm ra phía ngoài tế bào. Độ lớn của lực điện bằng: $F= qE = 3,2\cdot10^{-19}.8,75\cdot10^{6} = 28\cdot$ 10^{-13} N.
}
\end{ex}

% ===== Câu 6 =====
% ID: L11C3B18-03-TN
% Mức: NB
\begin{ex}
% ID: L11C3B18-03-TN
% Mức: NB
Trong ống phóng tia $X$ ở Bài Câu 4, một electron có điện tích $e = -1$, $6.10\,\mathrm{C}$ bật ra khỏi bản cực âm (catôt) bay vào điện trường giữa hai bản cực. Lực điện tác dụng lên electron đó bằng $\nguon{ly11 kntt.pdf}$
\choice
{\True $8 \cdot  10^{-13} N.$}
{$8 \cdot  10^{-18} N.$}
{$3,2 \cdot  10^{-17} N.$}
{$8 \cdot  10^{-15} N.$}
\loigiai{
Chọn A.

$F= qE = 1,6\cdot10^{-19}.5\cdot10^{6} = 8\cdot10^{-13}N$
}
\end{ex}

% ===== Câu 7 =====
% ID: L11C3B18-04-TL
% Mức: TH
\begin{ex}
% ID: L11C3B18-04-TL
% Mức: TH
Cho hai tấm kim loại phẳng rộng, đặt nằm ngang, song song với nhau và cách nhau $d =$ 5 $\mathrm{cm}$. Hiệu điện thế giữa hai tấm đó bằng $500\,\mathrm{V}$. a) Tính cường độ điện trường trong khoảng giữa hai bản phẳng. b) Khi một electron bật ra khỏi bản nhiễm điện âm và đi vào khoảng giữa hai bản phẳng với tốc độ ban đầu v_0 ≈ 0, hãy tính động năng của electron trước khi va chạm với bản nhiễm điện dương. $\nguon{ly11 kntt.pdf}$
\choiceTF

\loigiai{
a) Điện trường ở giữa hai tấm kim loại là điện trường đều, các đường sức song song và cách đều nhau và có cường độ điện trường bằng: $E = \dfrac{U}{d} = \dfrac{500}{0,05} = 10000\text{V}/\text{m}$

b) Động năng của electron trước khi va chạm với bản nhiễm điện dương sẽ bằng công của lực điện trường cung cấp cho electron trong dịch chuyển từ bản nhiễm điện âm sang bản nhiễm điện dương

$W$ = $A$ = Fd = |q|Ed = 1,6$\cdot$ 10^{-19}$\cdot$ 10^{4}.0,05 = 8$\cdot$ 10^{-17}J.
}
\end{ex}

% ===== Câu 8 =====
% ID: L11C3B18-04-TN
% Mức: NB
\begin{ex}
% ID: L11C3B18-04-TN
% Mức: NB
Khi một điện tích chuyển động vào điện trường đều theo phương vuông góc với đường sức điện thì yếu tố nào sẽ luôn giữ không đổi? $\nguon{ly11 kntt.pdf}$
\choice
{\True Gia tốc của chuyển động.}
{Phương của chuyển động.}
{Tốc độ của chuyển động.}
{Độ dịch chuyển sau một đơn vị thời gian.}
\loigiai{
Chọn A.

Khi một điện tích chuyển động vào điện trường đều theo phương vuông góc với đường sức điện thì gia tốc của chuyển động sẽ luôn giữ không đổi.
}
\end{ex}

% ===== Câu 9 =====
% ID: L11C3B18-05-TL
% Mức: TH
\begin{ex}
% ID: L11C3B18-05-TL
% Mức: TH
Hãy tính vận tốc theo phương Oy và động năng của electron khi va chạm với bản phẳng nhiễm điện dương ở câu 17* $\nguon{ly11 kntt.pdf}$
\choiceTF

\loigiai{
Vận tốc theo phương Oy của electron trước lúc va chạm với bản nhiễm điện dương là: $v_{y} = at = 3,516 \cdot 10^{13}\dfrac{x_{\max}}{20000} = 2,054 \cdot 10^{6}\left( \text{m/s} \right)$

Động năng của electron trước khi va chạm với bản nhiễm điện dương $(v_(x) = 0)$ là:

$W = \dfrac{1}{2}mv^{2} = \dfrac{1}{2}m\left( {v_{y}^{2} + v_{x}^{2}} \right) = \dfrac{1}{2} \cdot 9,1 \cdot 10^{-31} \cdot 4,22 \cdot 10^{12} = 19 \cdot 2 \cdot 10^{-19}\text{J}.$
}
\end{ex}

% ===== Câu 10 =====
% ID: L11C3B18-05-TN
% Mức: NB
\begin{ex}
% ID: L11C3B18-05-TN
% Mức: NB
Khi một điện tích chuyển động vào điện trường đều theo phương vuông góc với đường sức điện thì điện trường sẽ không ảnh hưởng tới $\nguon{ly11 kntt.pdf}$
\choice
{gia tốc của chuyển động.}
{\True thành phần vận tốc theo phương vuông góc với đường sức điện.}
{thành phần vận tốc theo phương song song với đường sức điện.}
{quỹ đạo của chuyển động.}
\loigiai{
Chọn B.

Khi một điện tích chuyển động vào điện trường đều theo phương vuông góc với đường sức điện thì điện trường sẽ không ảnh hưởng tới thành phần vận tốc theo phương vuông góc với đường sức điện.
}
\end{ex}

% ===== Câu 11 =====
% ID: L11C3B18-06-TN
% Mức: NB
\begin{ex}
% ID: L11C3B18-06-TN
% Mức: NB
Quỹ đạo chuyển động của một điện tích điểm q bay vào một điện trường đều $E$ theo phương vuông góc với đường sức không phụ thuộc vào yếu tố nào sau đây? $\nguon{ly11 kntt.pdf}$
\choice
{Độ lớn của điện tích q.}
{Cường độ điện trường E.}
{\True Vị trí của điện tích q bắt đầu bay vào điện trường.}
{Khối lượng m của điện tích.}
\loigiai{
Chọn C.

Quỹ đạo chuyển động của một điện tích điểm q bay vào một điện trường đều $\overset{\rightarrow}{E}$ theo phương vuông góc với đường sức không phụ thuộc vào vị trí của điện tích q bắt đầu bay vào điện trường.
}
\end{ex}
